% !TEX root = ../main.tex
\lecture{8}{2026-11-09}{Leverage Effekt und Banking}

\section{Leverage Effekt}

\begin{itemize}
    \item Abhängigkeit der EK-Rendite vom Verhältnis FK/EK. Untervoraussetzung: 
    \begin{itemize}
        \item Gesamtkapitalrendite bleibt konstant
        \item Keinen Einfluss auf Operative Geschäft
        \item FK-Zins konstant
    \end{itemize}
    \item Reingewinn = $ r_ek * EK$
\end{itemize}

\[
    r_{ek} = r_{gk} + \frac{FK}{EK}(r_{gk} - k_{fk})
\]

\begin{itemize}
    \item Hebel
    \begin{itemize}
        \item $\frac{FK}{EK}$, je grösser FK in Verhältnis zu EK, desto grösser die Rendite
        \item jedoch nur während die Gesamtkapitalrendite grösser ist als die FK-kosten
    \end{itemize}
\end{itemize}

\subsection{Rechenbeispiele}

\begin{itemize}
    \item 100\% EK = 300
    \begin{itemize}
        \item Gewinn Brutto => Ertrag - Aufwand = 30
        \item FK-Zinsen = 0 
        \item Reingewinn nach Zinsen = 30
        \item $r_ek = 30/300 = 10\% = r_gk $
    \end{itemize}
     \item 50\% FK = 150
     \begin{itemize}
         \item Gewinn Brutto => Ertrag - Aufwand = 30
         \item FK-Zinsen = 6\% (Gegeben) * 150 = 9
         \item Reingewinn nach Zinsen = 30 -9 = 21
         \item$ r_ek = 21/150 = 14\%  $
         \item Hebel hat positiver Effekt weil$ r_gk > r_fk  10\% > 6\%$
     \end{itemize}
       \item 67\% FK = 300/100 * 67 = 200
     \begin{itemize}
         \item Hebel hat auch positiver Effekt weil $r_gk > r_fk $
     \end{itemize}
\end{itemize}

\subsubsection{Fallunterscheidung}

Gewinn Schlechter/Guter/Optimaler Fall: 9 / 30 / 51

\begin{itemize}
    \item 100\% EK = 300
    \begin{itemize}
        \item $3\% / 10\% / 17\%$
        \item Avg 10\%
        \item Streuung vom Durchschnitt +/- 7\%
    \end{itemize}
    \item 50\% = 150
    \begin{itemize}
        \item FK Zinsen = 9 
        \item Reingewinn 0 / 21 / 42
        \item $r_ek = 0\% / 14\% / 28\%$
        \item $Avg 14\%$
        \item Streuung vom Durchschnitt +/- 14\%
    \end{itemize}
    \item 67\% Fk
    \begin{itemize}
         \item FK Zinsen = 12 
        \item Reingewinn -3 / 21 / 42
        \item $r_ek = 0\% / 14\% / 28\%$
        \item Avg 14\%
        \item Streuung vom Durchschnitt +/- 21\%
    \end{itemize}
    \item Insgesamt ergeben sich mit steigendem FK zu EK höhere durchschnittliche Renditen aber grössere Schwankungen/Risiko, dass die erzielte Rendite beim schlechten/mittleren Fall sogar negativ sein kann ($r_gk < r_ek$)
\end{itemize}

\section{Intro Banking}
Grundsätzlich besteht die primäre lukrative Tätigkeit von Banken aus dem langfristigen Finanzieren von Hypotheken, Unternehmenskredite aus kurzfristigen Anlagen von Privatpersonen.
\subsection{Grundproblematiken}
\begin{itemize}
    \item Informationsproblem
    \begin{itemize}
        \item Vermittlung von Kapitalgeber und -nehmer (Umverteilung des Kapitals der Gesellschaft, um Kapitalumlauf zu sichern)
    \end{itemize}
    \item Fristigkeitsproblem
    \begin{itemize}
        \item Langfristige Kreditvergabe mit kurzfristigen Anlagen
    \end{itemize}
    \item Losengrössenproblem
    \begin{itemize}
        \item Viele kleine Bezüge zu grossen Kapitalblöcke (Bankrun)
    \end{itemize}
    \item Risikoproblem
    \begin{itemize}
        \item Private Anleger möchten Sicherheit während die Bank das Geld unter Risiken einsetzen muss 
    \end{itemize}
\end{itemize}

\subsection{Funktionen von Banken}
Banken sind Finanzintermediäre die folgendes betreiben:
\begin{itemize}
    \item Kreditgeschäft
    \begin{itemize}
        \item Losgrössentransformationen (Aus kleinen -> grossen Kapitalmengen)
        \item Fristentransformation 
        \item Risikentransformation 
    \end{itemize}
    \item Verwahrung und Anlage Wertgegenstände
    \item Geldwechsel und Zahlungsverkehr (Operationen)
    \item Weitere interne Aufgaben:
    \begin{itemize}
        \item Risikomanagement 
        \item Monitoring
        \item Life-cycle Capital Allocation (Optimale Verteilung von Kapital zwischen Privatpersonen und )
        \item Transmission der Geldpolitik
        \begin{itemize}
            \item Zentralbanken bestimmen Geldpolitische Massnahmen (Zinssätze) die über die Privatbanken vermittelt werden 
        \end{itemize}
    \end{itemize}
\end{itemize} 
\footnote{"Weil Bank Risiken eingeht, muss sie entschädigt werden" reine Apologetik, es wird nicht versucht den Mechanismus zu erklären, der es möglich macht "endlos" aus Wert mehr Wert zu machen.}

\subsection{Akteure des Schweizerischen Finanzplatzes: SNB}

\subsubsection{Funktionen}
\begin{itemize}
    \item Liquiditätsversorgung
    \item Bargeldversorgung
\item Finanzstabilität (Inflation, Krisenverhinderung)
\item Verwaltung Währungsreserven
\end{itemize}
\subsubsection{Instrumente}
\begin{itemize}
    \item Festlegung des Leitzinses
    \item Offenmarktoperationen zur Liquiditätssteuerung
    \item Verzinzung 
    \item Devisenmarktinterventionen
\end{itemize}

\subsubsection{Repurchase Agreement (Repo)}
\begin{itemize}
    \item Primäres Instrument zur Steuerung der Geldmenge
    \item Wertpapiere (Repo) werden gegen einen Zinssatz durch SNB ausgestellt
\end{itemize}

